\section{Main algorithm and model-specific guarantees}
\label{sec:main-guarantee}

The estimand is the population target effect $\tau_t=\E_t\{Y(1)-Y(0)\}$.
There are $K$ sources and a target sample of size $n_t$. A common, prespecified
feature map $\phi$ includes an intercept. The target and source outcomes lie
in $[0,1]$. Conditional on all covariates and treatment assignments
$\mathcal D$, patient outcomes are independent. Target patients have i.i.d.
covariates from the target population. Identification of the potential-outcome
means requires the usual consistency, conditional exchangeability and overlap
conditions; the statistical model below concerns those identified arm means.

The target arm means are $m_t^a(x)=\phi(x)^T\beta_t^a$. For each arm, a fixed
but unknown set of at least $g_a\ge1$ sources has this same conditional mean.
The two valid sets need not coincide. The numbers $g_a$ are assumptions, not
counts estimated by comparing outcomes. Nonzero source deviations have no
minimum magnitude. We distinguish three source model classes:
\begin{itemize}
\item \textbf{Linear.} Every source arm mean, including invalid ones, belongs to the
common linear span, with unrestricted site-specific coefficients.
\item \textbf{Declared approximation.} Target and valid-source means remain linear.
Invalid means satisfy a declared design-conditional projection bound
$\|(I-H_{ja})m_{ja}\|_2\le\rho_{ja}$. An estimated certificate instead
requires a simultaneous failure allowance $\delta_{\rm app}$.
\item \textbf{Bounded invalid.} Only target and valid-source means must be linear;
the other means can be arbitrary in $[0,1]$. Set
$\rho_{ja}=\sqrt{n_{ja}}/2$ for the centered mean vector and pay no additional
certificate failure probability.
\end{itemize}
Approximate valid or target means require the extra mean-bias allowances of
Section~\ref{sec:approximation-control}. They are not covered by the unmodified
linear-target coefficient certificate below.

\subsection{One executable algorithm}

Fix the model class, valid counts, feature map, design policy and error ledger
before examining outcomes. The following steps define the current procedure.
\begin{enumerate}
\item Compute $b_t=\mathbb P_t\phi(X)$. Check each arm's rank, condition number
and leverage using only $\mathcal D$. A target failure returns $[-1,1]$.
Under \texttt{all\_required}, any source failure uses the protected target
procedure. Under \texttt{remove}, drop every site failing either arm, let
$m$ be the number dropped, and set
\[
 K'=K-m,\qquad g_a'=\max(0,g_a-m).
\]
If $K'=0$ or either $g_a'=0$, use the protected target procedure. Otherwise
perform every subsequent source calculation on the retained set, with $K'$,
$g_a'$ and $g_\cap'=\max(0,g_1'+g_0'-K')$.
\item At each retained source and the target, solve
$X_{ja}^T\gamma_{ja}=b_t$ with minimum Euclidean norm. Return
$Z_j^a=\gamma_{ja}^TY_{ja}$, the leave-out variance estimate
$\widehat v_{ja}$, and the design summaries in
Section~\ref{sec:continuous-balance}. Do not truncate a negative
$\widehat v_{ja}$ before dispersion correction.
\item For each arm, form the variance-corrected source dispersion using
$K'$. In the approximation modes, at most $q_a'=K'-g_a'$ retained sites can
have nonzero projection errors. Recompute the top-$q_a'$ bounds $C_B,C_U$,
the shifted statistic $\widehat D+a_{K'}C_B$, and the corrected MGF constant
$Q_0^{\rm app}$. For a contrast, add the two arm-specific approximation
budgets. The inversion in Sections~\ref{sec:continuous-balance}--\ref{sec:approximation-control}
gives simultaneous upper bounds $U_1,U_0,U_C$.
\item With $g_\cap'>0$, the adaptive dispersion radius is
\[
 B=\min\left\{
 \sum_{a=0}^1\sqrt{\frac{K'-g_a'}{g_a'}U_a},\quad
 \sqrt{\frac{K'-g_\cap'}{g_\cap'}U_C}\right\}.
\]
The second option is omitted when $g_\cap'=0$. Allocate $\alpha_D/4$ to
each arm and $\alpha_D/2$ to the contrast when both options are used;
otherwise allocate $\alpha_D/2$ per arm. A guaranteed all-valid arm has zero
bias contribution. The source and target bands are
\begin{align*}
 J_S&=\left[\frac1{K'}\sum_j(Z_j^1-Z_j^0)
       \ \pm\left\{\sqrt{\frac{\log(2/\alpha_S)}{2(K')^2}
                         \sum_{ja}\|\gamma_{ja}\|_2^2}+B\right\}\right],\\
 J_T&=\left[Z_t^1-Z_t^0\ \pm
          \sqrt{\frac{\log(2/\alpha_T)}2\sum_a\|\gamma_{ta}\|_2^2}\right].
\end{align*}
Use $J_S\cap J_T$ if nonempty, and the shorter band otherwise. A source
design/count fallback uses $J_T$ directly. Denote this conditional interval
by $I_X$; it targets the observed-target effect
$\tau_X=b_t^T(\beta_t^1-\beta_t^0)$.
\item Construct the target slope-contrast ellipsoid with failure
$\alpha_R$. Its simultaneous upper bounds $R^+,V_X^+$ on the true CATE
range and empirical variance give
\[
 r_X=\sqrt{\frac{2V_X^+\log(4/\alpha_X)}{n_t}}
       +\frac{7R^+\log(4/\alpha_X)}{3(n_t-1)}.
\]
Expand $I_X$ by $r_X$ and clip to $[-1,1]$. First project the target point onto
$I_X$, then clip that projected value to the reported population interval to
define the point estimate. No simulated true CATE
range, variance or source labels are supplied to this construction.
\end{enumerate}

The implemented default ledger is
$(\alpha_S,\alpha_T,\alpha_D,\alpha_R,\alpha_X)
=(.01,.01,.015,.005,.01)$. A separately reported protected target-only
comparator uses $(\alpha_T,\alpha_R,\alpha_X)=(.0225,.005,.0225)$.
Thus it benefits from the same target certificate and its own full $.05$
budget. The ordinary normal target interval is an asymptotic comparator.

\begin{theorem}[Coverage, removal and the scope of precision improvement]
\label{thm:main}
Under any of the three declared source classes and the target conditions
above, the algorithm satisfies, for every finite $K$,
\[
 \Prb\{\tau_t\in I\}\ge
 1-\alpha_S-\alpha_T-\alpha_D-\alpha_R-\alpha_X-\delta_{\rm app}.
\]
Here $\delta_{\rm app}=0$ for deterministic valid approximation envelopes
and the universal bounded-invalid mode. Conditional on the design and valid
deterministic approximation bounds, the conditional interval has coverage at
least $1-\alpha_S-\alpha_T-\alpha_D$ for $\tau_X$. Source/design fallback
requires only its target event; target-design failure covers automatically.
These claims need no separation of invalid-source biases from zero, and hold
under either design policy.

For precision, suppose on a design event $\mathcal G_n$ that $K'\asymp K$,
$g_a'/K'\ge c_0>0$, balanced arm sizes are of order $n_s$, and the retained
design factors obey $S_j=O(n_s^{-1})$, $G_j,F_j=O(n_s^{-3})$ and
$d_j^{\max}=O(n_s^{-2})$, uniformly. In the linear class, if both arm
dispersions are $O(n_s^{-1}K^{-1/2})$, the expected conditional source width
on this event is $O(n_s^{-1/2}K^{-1/4})$. A positive contrast intersection
is not necessary when arm certificates are used.

With fixed target dimension and the well-conditioned target designs of
Section~\ref{sec:target-protection}, put
\[
 r_n=n_s^{-1/2}K^{-1/4}
       +\sigma_\tau n_t^{-1/2}+n_t^{-1}.
\]
The unconditional expected reported length is bounded by
$C r_n+2\Prb(\mathcal G_n^c)$. Obtaining order $r_n$ therefore also requires
$\Prb(\mathcal G_n^c)=O(r_n)$. For growing target dimension use the universal
$n_t^{-1/2}$ composition bound unless the stronger target rate has separately
been established.

For the approximation classes, let $A_B$ uniformly bound the relevant
$a_{K'}C_B$, and let $A_Q$ bound $Q_0^{\rm app}-Q_0$. A sufficient modified
width order is $r_n+\sqrt{A_B}+A_Q^{1/4}$, with the same fallback contribution.
Consequently, bounded-invalid coverage alone does not imply the linear
fourth-root improvement.
\end{theorem}

\begin{proof}
Given $\mathcal D$, removal fixes a retained set without observing outcomes.
Deleting $m$ coordinates removes at most $m$ members of each valid set,
proving $g_a'$; hence at most $K'-g_a'$ retained arm means need the projection
correction. The corrected quadratic upper bounds and Hoeffding source/target
bands hold simultaneously by their error allocation. The valid-count
dispersion inequality then covers $\tau_X$ in $J_S$, while $J_T$ covers it
on its target event. On their joint event the intersection is nonempty and
covers; changing the empty-intersection branch cannot remove this event.
The count/design fallback uses the same target bound.

The target coefficient event dominates both true CATE range and empirical
variance. Apply empirical Bernstein to the fixed true CATE function, then
union this event with the coefficient and conditional outcome events. No
independence between those events is required. This also accommodates
design-dependent source removal without conditioning away the i.i.d. target
composition argument. Add $\delta_{\rm app}$ if an envelope is estimated.

For length, the conditional expected-dispersion calculation in
Section~\ref{sec:design-stability} gives the stated linear rate. Under
approximation, $|\E\widehat D-D|\le a_{K'}C_B$ and
$\Var(\widehat D\mid\mathcal D)\le LD+Q_0^{\rm app}$. The same inversion,
triangle inequality and Jensen argument yield an additional
$O(\sqrt{A_B}+A_Q^{1/4})$ radius. Intersection or shorter fallback is no
longer than the source band on $\mathcal G_n$. Add the target composition
expectation bound. Outside $\mathcal G_n$ the clipped interval has length
at most two. This establishes the expected-length statement.
\end{proof}

The theorem is uniform over arbitrary weak source deviations within each
declared class for coverage. Its length assertion is an adaptation result on
the stated small-dispersion and stable-design subclass; it does not promise
short intervals at every fixed, strongly heterogeneous configuration.
